\documentclass[12pt,letterpaper]{amsart}
\usepackage[T1]{fontenc}
\IfFileExists{newtxtext.sty}%
  {\usepackage{newtxtext,newtxmath}}%
  {\usepackage{mathptmx}\usepackage{amssymb}\renewcommand{\ttdefault}{lmtt}}
\usepackage[letterpaper,margin=1in]{geometry}
\usepackage{microtype}
\usepackage{booktabs,longtable,array,enumitem}
\usepackage[hidelinks]{hyperref}
\emergencystretch=2em

\title[Draft 3 Revision Log]{Revision Log for\\\emph{On a metaplectic Jacquet--Langlands transfer\\for division algebras of prime degree}\\[0.5em]\normalsize Draft 3 (Full Book Build)}
\author{David Joyner}
\address{\texttt{wdjoyner@gmail.com}}
\date{November 15, 2026}

\begin{document}
\maketitle

\section{Overview and provenance}
This file records every mathematical correction, verification step, style-guide adjustment, and build fix made in passing from Draft~2 (November~12, 2026) to Draft~3 (November~15, 2026) of the full monograph (\path{mjl_book_draft3.tex} and the thirteen included files in \path{chapters/}).

The five core chapters of the monograph originated as standalone manuscripts in \url{https://github.com/wdjoyner/metaplectic-jl}:
\begin{itemize}[leftmargin=*]
\item Chapter~4 (\path{chapters/ch4_part1.tex}): Part~I, Draft~11;
\item Chapter~5 (\path{chapters/ch5_part2.tex}): Part~II, Draft~10;
\item Chapter~6 (\path{chapters/ch6_part3.tex}): Part~III, Draft~10;
\item Chapter~7 (\path{chapters/ch7_part4.tex}): Part~IV, Draft~6;
\item Chapter~8 (\path{chapters/ch8_wild.tex}): Companion wild rank-two paper, Draft~8.
\end{itemize}
No proof status tag (\textbf{(P)}, \textbf{(E)}, \textbf{(O)}, \textbf{(C)}, or \textbf{(P$^{\mathrm c}$)}) was upgraded silently in Draft~3. All 495 original \verb|\label| identifiers and theorem numbers across Chapters~4--8 are preserved unchanged.

\section{Responses to the November 14, 2026 referee report}
Each point in ChatGPT's November~14, 2026 referee report was checked independently before any edit was applied (Style Guide \S\S4.5, 7.4).

\begin{enumerate}[leftmargin=*,label=\textbf{R\arabic*.}]
\item \textbf{Response to E1 (Local division algebras of degree $m>2$, Chapter~3).} File: \path{chapters/ch3_notation.tex}.
\emph{Independent check (Level 1):} Over a nonarchimedean local field $F$, Hasse's theorem gives a canonical isomorphism $\operatorname{Br}(F)\xrightarrow{\sim}\mathbb Q/\mathbb Z$, under which central division algebras of degree $m\ge 1$ correspond bijectively to elements $a/m\bmod 1$ of exact order $m$, i.e., $1\le a\le m$ with $\gcd(a,m)=1$. There are $\varphi(m)$ non-isomorphic central division algebras of degree $m$ over $F$, realized as cyclic algebras $D_{a/m}=F_m\langle\Pi\rangle$ with $\Pi^m=\varpi$ and $\Pi x\Pi^{-1}=\sigma^a(x)$. Uniqueness up to isomorphism holds if and only if $\varphi(m)=1$, namely $m\in\{1,2\}$. Draft~2's assertion in \path{ch3_notation.tex} that there is a unique local division algebra for every $m\ge 1$ was false for $m\ge 3$ and contradicted Chapter~8 (\path{ch8_wild.tex}, \S8.2.1).
\emph{Action taken:} Replaced the ``Local classification; the standard construction'' item in \path{chapters/ch3_notation.tex} (\S3.3) with the exact $\varphi(m)$ classification and the $\sigma^a$ cyclic presentation, citing \texttt{Ser79} and \texttt{Rei75}, and updated the ``Embedding criterion'' wording in \S3.3 from ``the degree-$m$ division algebra'' to ``every degree-$m$ division algebra over $F$''.

\item \textbf{Response to E2 (Cocycle existence vs.\ scalar centrality, Chapter~1).} File: \path{mjl_book_draft3.tex}.
\emph{Independent check (Level 1):} Because $\sigma_{K_v}(g,h)$ and $(\det g,\det h)_{n,K_v}$ are both $2$-cocycles on $\mathrm{GL}(r,K_v)$ valued in $\mu_n$, the product $\sigma_{K_v}(g,h)(\det g,\det h)_{n,K_v}^{c_v}$ satisfies the $2$-cocycle identity for \emph{every} $c_v\in\mathbb Z$. On scalar matrices $aI_r,bI_r$, the cocycle evaluates to $(a,b)_{n,K_v}^{e_r}$ with $2e_r = r(r-1+2rc_v)$ (as already computed in \path{ch6_part3.tex}, equation~(6.6.2)), and diagonal commutators against scalars have exponent $r-1+2rc_v$ (\path{ch7_part4.tex}, Proposition~7.2.1). For odd $n$ coprime to $r$, $n\mid (r-1+2rc_v)$ is therefore the exact condition for scalar centrality, not for cocycle existence.
\emph{Action taken:} Corrected \S1.1 (``The two covers'') in \path{mjl_book_draft3.tex} and added worked Example~1.1.1 contrasting $c_v=0$ and $c_v=2$ for $n=3,r=2$.

\item \textbf{Response to E3 (Scope of wild integrality in Chapter~8).} Files: \path{mjl_book_draft3.tex}, \path{chapters/ch6_part3.tex}, and \path{chapters/ch7_part4.tex}.
\emph{Independent check (Level 2 against Chapter 8 and Chapter 9):} Chapter~8 proves Target~A ($L_D\mathbf 1$ is genuine irreducible of dimension $27$, Theorem~8.11.1) and Target~B for representations trivial on $K^3$ and their norm twists (levels $0$, $1$, and $2$, Theorems~8.12.5 and~8.12.7) modulo exhaustive exact computation (\textbf{(P$^{\mathrm c}$)}), while explicitly leaving congruence levels $\ge 3$ open (\S8.14).
\emph{Action taken:} Removed ``completely'' from both occurrences in Chapter~1 (\S1.2 and \S1.3), stated the exact scope (trivial character and representations trivial on $K^3$, i.e., through level two, modulo \textbf{(P$^{\mathrm c}$)}) in the Abstract, \S1.2, \S1.3, and \S1.4, added worked Example~1.4.1 on level~$3$, and harmonized the running-example forward pointers in \path{ch6_part3.tex} (\S6.1) and \path{ch7_part4.tex} (\S7.1) so they cross-reference Chapter~8's partial resolution consistently.

\item \textbf{Response to E4 (Audit of \textnormal{\textbf{(P$^{\mathrm c}$)}} tables and build).}
\emph{Independent check (Level 1):} Wrote and executed an independent pure-Python exact modular verification script (\path{scripts/verify_appendix_tables.py}) checking all Target~A finite-design tables in Appendix~A (\path{chapters/appA_scripts.tex}):
\begin{itemize}[leftmargin=*]
\item Verified that $\operatorname{Ad}(1+j)$ and $\operatorname{Ad}(1+ij)$ generate a group of order $27$, $\operatorname{Ad}(1+i)$ and $\operatorname{Ad}(j)$ generate $D_8$ of order $8$, and together they generate a group of order $216$ on $V=\mathbb F_3^5$.
\item Verified that the $36$ good lines ($72$ vectors of $\bar G_3$) are invariant under the $216$-element group, and that all $B=720$ ordered coplanar triples of distinct good lines have zero defect sum modulo~$3$.
\item Verified across all $3^5=243$ covectors $h\in (\mathbb F_3^5)^\vee$ that $N_{e_*}(h)$ has exact distribution $\{11:216,\ 20:27\}$, and that the $27$ covectors achieving $20$ match the printed set $\Lambda=\Lambda_+$ identically.
\item Verified that the cocycle covectors $\beta_{\operatorname{Ad}(1+j)}=(1,0,0,1,0)$ and $\beta_{\operatorname{Ad}(1+ij)}=(0,1,0,0,1)$ satisfy $e_*(Mv)-e_*(v)\equiv \beta_M\cdot v\pmod 3$ for all $72$ vectors $v\in\bar G_3$.
\end{itemize}
All \textbf{(P$^{\mathrm c}$)} labels are retained unchanged pending independent human review and full re-execution of the repository's Hilbert-symbol scripts (\path{part-IV/scripts/}).
\end{enumerate}

\section{Additional mathematical, bibliographic, and build fixes}
\begin{enumerate}[leftmargin=*]
\item \textbf{Correction of $b_D$ and $c_D$ descriptions in Example 1.3.1 (\path{mjl_book_draft3.tex}).}
Draft~2's parenthetical gloss in Example~1.3.1 called $b_D=3$ ``the forced common dimension of genuine representations'' and $c_D=27$ ``the discriminant of the relevant congruence filtration''. By Theorem~6.4.3, Corollary~6.5.2, and Proposition~6.7.1 in \path{ch6_part3.tex}, $b_D=|n|_F^{-(r-1)/2}=3$ is the unitary normalization factor of $L_D$ (and the forced common divisor of genuine irreducible dimensions from covered elliptic tori), whereas $c_D=|n|_F^{-(r^2-1)/2}=27$ is the near-center dimension-scaling factor ($\dim\tau = c_D\dim\rho$). Updated Example~1.3.1 accordingly.
\item \textbf{Font fallback fix under \texttt{pdflatex} (\path{mjl_book_draft3.tex}).}
Set \texttt{\textbackslash ttdefault} to \texttt{lmtt} inside the \texttt{mathptmx} fallback branch so that \texttt{T1}-encoded typewriter text uses Type~1 Latin Modern Typewriter rather than missing \texttt{ectt1200} bitmap fonts when \texttt{newtxtext.sty} is not installed.
\item \textbf{Elimination of Overfull \texttt{\textbackslash hbox} and \texttt{hyperref} warnings.}
Replaced \verb|\texttt| with breakable \verb|\path| on the three long script filenames in \path{chapters/appA_scripts.tex} (item~7), removing the $11.15\,\mathrm{pt}$ overfull box without manual line breaks (Style Guide \S1.3). Added \verb|\texorpdfstring| to \verb|\section| headings in \path{chapters/ch5_part2.tex} (\S5.6) and \path{chapters/ch8_wild.tex} (\S8.8), eliminating all five \texttt{hyperref} PDF-bookmark warnings.
\item \textbf{Elimination of hand-typed bracket citations and typos (Style Guide \S1.5).}
Replaced \texttt{[DDMS; SW]} in \path{chapters/ch2_literature.tex} (\S2.4) with \verb|\cite{SW00}| alongside \verb|\cite{Laz65}|; replaced \texttt{[Serre, Local Fields, Ch.~XIV]} in \path{chapters/ch3_notation.tex} (\S3.5) with \verb|\cite[Ch.~XIV]{Ser79}|; fixed doubled ``\texttt{using (using (E))}'' and straight ASCII quotation marks in \path{chapters/ch3_notation.tex} (\S3.1).
\item \textbf{Exposition and worked examples in Chapter 1 (Style Guide \S3).}
Removed all prose semicolons in Chapter~1, added orienting paragraphs to \S\S1.1--1.4, and added worked Examples~1.1.1, 1.2.1, and~1.4.1 without altering the numbering of Example~1.3.1 (\verb|ex:running|) or Remark~1.3.2 (\verb|rem:running-hyp|).
\end{enumerate}

\section{Default rulings used on \texttt{[DECIDE]} items and external sources}
\begin{itemize}[leftmargin=*]
\item \textbf{Document class and Times font:} Kept the \texttt{book} class (\texttt{12pt, letterpaper, oneside}) and \texttt{newtxtext}/\texttt{newtxmath} with the \texttt{mathptmx} and \texttt{lmtt} fallback.
\item \textbf{Glossary placement:} Kept Chapter~3 (Background and notation) near the front and Chapter~10 (Glossary) near the back before the Acknowledgments.
\item \textbf{Author line:} Kept the Draft~2 title-page author list (\emph{David Joyner, Claude (Anthropic), ChatGPT (OpenAI)}) unchanged and recorded Gemini (Google)'s Draft~3 revision and verification role in the Abstract, Acknowledgments, and this Revision Log.
\item \textbf{Sources not re-consulted line by line (Style Guide \S4.3):} As recorded in \path{ch2_literature.tex}, Klose's papers \texttt{Klo86a}, \texttt{Klo86b} and dissertation were not directly consulted; theorem-by-theorem re-verification of external \textbf{(E)} citations in Chapters~4--8 against physical journal pages was not performed during this pass.
\item \textbf{Build audit:} Three passes of \texttt{pdflatex} plus \texttt{makeindex} on \path{mjl_book_draft3.tex} produced \path{mjl_book_draft3.pdf} (216 pages) with $0$ errors, $0$ undefined references or citations, $0$ overfull boxes, and $0$ package warnings.
\end{itemize}

\section{Corrections after review of Draft 3 (Claude, November 15, 2026)}
Claude reviewed Gemini's Draft~3 against Draft~2 and the November~14 referee report, rebuilt the book (216 pages, no undefined references, no overfull boxes), and reran \path{scripts/verify_appendix_tables.py} (all checks pass). That script checks that the printed Appendix~A tables are consistent with one another. It does not derive the adjoint matrices or $e_*$ from the quaternion arithmetic or the Hilbert symbol, so the \textbf{(P$^{\mathrm c}$)} labels stay as they are. The following corrections were then made. No status label was changed.
\begin{enumerate}[leftmargin=*,label=\textbf{C\arabic*.}]
\item \textbf{Elliptic classes (\S1.1, Motivation).} Draft~3 described the elliptic classes of $\mathrm{GL}(r)$ as those whose characteristic polynomial ``does not split''. That is wrong for $r\ge3$: a linear factor times an irreducible quadratic does not split but is not elliptic. The sentence now says that regular elements of $D^\times$ correspond to regular elliptic classes, whose characteristic polynomials are irreducible. In the same paragraph, ``acts faithfully by a fixed character'' became ``acts through a fixed faithful character''.
\item \textbf{Absolute value in Example~1.3.1.} Draft~3 introduced $|\cdot|_p$ with $|p|_p=p^{-1}$ and then used $|3|_{K_{\mathfrak p_3}}=3^{-2}$, which is the normalized absolute value of $\mathbb Q_3(\zeta_3)$. The example now defines the normalized absolute value $|\cdot|_F$ of a local field and says why $|3|=3^{-2}$ there.
\item \textbf{Order of examples in Chapter~1.} Gemini's Example~1.2.1 (the trivial representation traced through the chapters) used the running example before Example~1.3.1 defines it. It now comes at the end of \S1.3, as Example~1.3.3. Examples~1.3.1, 1.4.1 and Remark~1.3.2 keep their numbers.
\item \textbf{Terminology in Chapter~3.} After the local classification was rewritten as the ``cyclic construction'' $D_{a/m}$, three later items still referred to ``the standard construction''. They now refer to the cyclic construction, and the Cyclic algebras item notes that $D_{a/m}$ uses the generator $\sigma^a$.
\item \textbf{Scope of toral positivity.} The abstract and \S1.2 correctly restrict the quadratic toral permutation-positivity theorem of Chapter~7 to odd residual characteristic. Two other summaries of it omitted that hypothesis: \S7.6 (the subsection ``What the local reductions prove, and what they do not'') and the opening paragraph of \S8.5. Both now state it, together with $\mu_n\subset F$.
\item \textbf{Authorship and credit.} Ruling by David Joyner (November~15, 2026): Gemini's contribution to Draft~3 does not amount to authorship. The title page is unchanged. The abstract now says Claude and ChatGPT co-wrote the book and that Gemini prepared the Draft~3 revision. The Acknowledgments now credit Gemini in a paragraph of its own, describing its work in detail, and no longer say that Claude and ChatGPT are authors only of ``the underlying manuscripts''.
\end{enumerate}

\section{Referee requests deferred to a later draft}
The following requests in the November~14 report were not addressed in Draft~3: an audit table of the imported theorems with exact theorem and page citations; machine-readable logs, with source commit and dependency versions, for every load-bearing computation; further worked examples in Chapters~4--8; and shorter sentences in Chapters~4, 5 and~8, which still contain many long semicolon-joined sentences.

\end{document}
